\chapter{Mot Grunnlaust}

\begin{motto}
Eit syn som ikkje toler sitt eige sterkaste motsvar,\\ fortener ikkje å bli ståande. Her er motsvaret.
\end{motto}

No møter dei to bøkene kvarandre direkte. Eg tek dei berande påstandane i \textsc{Grunnlaust} éin for éin, i deira eiga form, og svarar på kvar med det sterkaste forsvaret har. Eg kjem ikkje til å vinne alle. Der \textsc{Grunnlaust} har rett, seier eg det — eit forsvar som hevdar å vinne alt, vinn ingenting. Men eg kjem til å vise at hovudslutninga — at faget manglar eit fundament og difor er ein avleggar — ikkje følgjer av premissane, sjølv når premissane er sanne.

\section{Om grunnsetninga}

\begin{innvending}
«Design kan ikkje grunngje dommane sine slik fysikken kan; difor manglar det eit fundament.»
\end{innvending}

\noindent Dette er kategorifeilen, og han er den djupaste feilen i heile \textsc{Grunnlaust}. Slutninga føreset at den einaste forma for fagleg fundament er den naturvitskaplege. Men det er ein føresetnad forkledd som ein konklusjon. Simon viste at design høyrer til vitskapane om det kunstige — om korleis ting \emph{kan} bli, ikkje korleis dei \emph{er} \autocite{simon1969sciences} — og Cross at design difor må drivast som ein disiplin på eigne premissar, ikkje som ein vitskap \autocite{cross2001discipline}. Jussen, medisinen, pedagogikken og historiefaget kan heller ikkje grunngje dommane sine slik fysikken kan, og likevel er dei modne fag. Anten fell alle desse med same slaget, eller så finst det meir enn éi form for fundament. \textsc{Grunnlaust} kan ikkje ha det begge vegar.

\section{Om den tause kunnskapen}

\begin{innvending}
«Den tause kunnskapen i atelieret er privat heile vegen ned; det er ikkje noko felles under han å formulere.»
\end{innvending}

\noindent Dette er empirisk feil. Harry Collins sin systematiske analyse viser at det meste av taus kunnskap er \emph{kollektiv} — samfunnets eigedom, tileigna gjennom sosialisering i ein praksisfellesskap — og difor delt, ikkje privat \autocite{collins2010tacit}. Polanyi viste at \emph{all} kunnskap, òg fysikken, kviler på ein taus dimensjon; om det å ha ein taus kjerne gjorde eit fag fundamentlaust, ville fysikken falle med design. Schön viste at den tause kunnskapen i designstudioet er eksaminerbar — coachbar, demonstrerbar, kritiserbar — altså intersubjektivt tilgjengeleg. At kunnskap er taus, tyder ikkje at han er privat. Det var det heile argumentet hang på, og det held ikkje.

\section{Om metoden og kumulasjonen}

\begin{innvending}
«Design har inga delt metode og ingen kumulativ kunnskap; forskinga akkumulerer ikkje.»
\end{innvending}

\noindent Her er \textsc{Grunnlaust} på sitt svakaste, fordi moteksempla er konkrete og talrike. Det finst heile, replikerbare metodologiar: Blessing og Chakrabarti si \textsc{drm} med valideringsfasar \autocite{blessing2009drm}. Det finst kumulativ kunnskap som til og med overførast \emph{mellom domene}: Alexander sitt mønsterspråk vart tilpassa til programvare av Gang of Four og er no fagstandard \autocite{gof1994patterns}. Det finst empirisk testa, replikerte designpåstandar: Ulrich sin kontrollerte studie i \textit{Science} av naturutsyn og liggjetid \autocite{ulrich1984view}. Og det nordiske miljøet har vist korleis design lagar si eiga teori gjennom praksis \autocite{redstrom2017making}. Påstanden om at ingenting akkumulerer, kan berre haldast oppe ved å omdefinere kvart moteksempel til «ikkje ekte kunnskap» — og det er nett den ufalsifiserbarheita \textsc{Grunnlaust} skuldar design for.

\section{Om det vrange problemet}

\begin{innvending}
«‘Vrange problem’ er eit alibi faget brukar for å sleppe å lage teori.»
\end{innvending}

\noindent Rittel, som mynta omgrepet, brukte det ikkje som alibi; han bygde rigorøse argumentasjonsmetodar som svar på vrangskapen \autocite{rittel1973}. Innsikta hans er ekte: designproblem \emph{er} strukturelt ulike tamme vitskapsproblem, og å krevje naturvitskapleg falsifiserbarheit av dei er ein kategorifeil. \textsc{Grunnlaust} har likevel eit poeng her, og eg gjev det: vrangskapen er spesifikasjonen på kva ein teori må handtere, ikkje ein grunn til å sleppe han. Men det følgjer ikkje at teorien er umogleg — berre at han må vere ein teori om navigasjon under vrangskap. At eit slikt prosjekt er vanskeleg, gjer det ikkje til eit alibi.

\section{Om migrasjonen av design thinking}

\begin{innvending}
«Design thinking migrerte til økonomifaget fordi design ikkje hadde nokon grunn å bli verande på.»
\end{innvending}

\noindent Den motsette lesinga er minst like sterk: design gjekk \emph{inn i} økonomien med noko økonomien ikkje sjølv kunne lage — meiningsinnovasjon \autocite{krippendorff2006}. Verganti viser at designens distinkte bidrag er å endre kva ting \emph{tyder}, ein tredje innovasjonsveg utover teknologi og marknad. Ein pode som slår rot i ny jord fordi han ber noko den jorda manglar, er ikkje rotlaus; han er fruktbar. \textsc{Grunnlaust} les migrasjonen som svakheit fordi boka alt har bestemt seg for at design er tomt. Les ein han utan den føresetnaden, ser ein ein disiplin som eksporterte ein reell kompetanse.

\section{Der Grunnlaust står}

Eg lova å seie kvar kritikken vinn, og her gjer han det. \textsc{Grunnlaust} hevdar at design manglar eit internasjonalt etisk organ med makt til å ekskludere, og at det ikkje finst nokon standard ein kan bli utelukka frå. Det er, med små atterhald, sant. Design har rike etikkar og kodeksar, men nesten ingen kan handhevast over individet slik arkitektane sine kan \autocite{riba2019code}. Faget har heller ikkje samla etikkane sine til éi felles, vedteken forplikting. Desse er reelle manglar, og forsvaret skal ikkje pynte på dei.

Men sjå kva slags manglar dei er. Dei er manglar ved \emph{institusjonell modning} — ved apparatet for å handheve og samle — ikkje ved intellektuelt fundament. Eit fag kan ha eit grunnlag og enno ikkje ha bygt tribunalet som vaktar det; det skildrar dei fleste profesjonar i deira fyrste hundreår. Skilnaden mellom \textsc{Grunnlaust} og \textsc{Grunnfest} kokar til slutt ned til dette: kritikaren les fråvêret av eit handhevingsapparat som prov på fråvêret av eit fundament, medan forsvaret les det som teikn på eit ungt fag som enno byggjer institusjonane sine. Den fyrste lesinga gjer design til ein avleggar. Den andre gjer det til ein disiplin under modning, som til og med har eit ontologisk og verdsskapande kall langt utover det kritikaren ser \autocite{escobar2018pluriverse}.

\section{Kva som blir ståande}

Når begge sider har sagt sitt, står dette att. \textsc{Grunnlaust} har rett i at design manglar eit handhevingsapparat og ein samla etikk, og rett i at faget har gjort reell skade når det optimaliserte éin akse blindt. \textsc{Grunnfest} har rett i at design har ei eiga kunnskapsform, ein metode, eit fagspråk og ein etikk — og at kravet om eit naturvitskapleg fundament er ein kategorifeil som ville felle halve universitetet om det vart brukt konsekvent. Begge kan ikkje ha rett om hovudslutninga. Men begge har rett om mykje, og det er det ærlege resultatet av å skrive begge bøkene: ikkje ein vinnar, men ei usemje som no er ført så skarpt som ho kan bli. Lesaren får avgjere. Det var alltid meininga.

\vignett

Det står att eit siste ord — ikkje eit argument, men ei vurdering av kva slags fag dette eigentleg er, og kva det kan bli. Det er etterordet.

\vidarelesing{\fullcite{cross2001discipline}; \fullcite{collins2010tacit}; \fullcite{gof1994patterns}; \fullcite{redstrom2017making}; \fullcite{rittel1973}; \fullcite{escobar2018pluriverse}.}
